% TOP_PROBLEM: {"id":30007535,"problem_number":"LOCAL-30007535","title":"Safe debt and fiscal limits","category_id":21,"category":{"id":21,"name":"economics","display_name":"Economics","description":"Open economic theory statements, published research agendas, and proposed scoped specifications; question provenance and historical priority are qualified per record.","slug":"economics","order_index":21,"created_at":"2026-10-08T00:00:00Z"},"status":"research_question","external_url":null,"legacy_ids":[],"published":false,"statement":"Identify a horizon- and institution-specific safe-debt capacity satisfying an explicit bond-loss probability constraint and fiscal feasibility.","economics":{"original_id":24,"field":"Fiscal policy and sovereign debt","kind":"design","formalization_basis":"adapted_research_question","source_status":"explicit_open","priority_status":"earliest_found","priority_note":"This is the earliest explicit relevant statement verified in this audit, not a claim of original priority; older literature and earlier versions may contain antecedents.","earliest_verified_reference":{"authors":["Kris James Mitchener","Christoph Trebesch"],"title":"Sovereign Debt in the 21st Century","year":2021,"url":"https://www.ifo.de/DocDL/cesifo1_wp8959.pdf","doi":"10.3386/w28598","locator":"Section 10, “Concluding Remarks,” pp. 51–52","evidence_summary":"The survey says many puzzles remain, including how much more sovereign debt markets can absorb and how official holdings, bailouts, panics, and sovereign–bank feedback affect future debt crises."},"current_reference":{"authors":["Kris James Mitchener","Christoph Trebesch"],"title":"Sovereign Debt in the 21st Century","year":2021,"url":"https://www.ifo.de/DocDL/cesifo1_wp8959.pdf","doi":"10.3386/w28598","locator":"Section 10, “Concluding Remarks,” pp. 51–52","evidence_summary":"The survey says many puzzles remain, including how much more sovereign debt markets can absorb and how official holdings, bailouts, panics, and sovereign–bank feedback affect future debt crises."},"formalization_note":"The survey explicitly asks how much debt markets can absorb. This probabilistic safety-capacity definition is an adapted operational target, not a published universal fiscal limit."},"statusQualification":"Published question or research agenda verified at the cited locator. The mathematical specification is adapted for this catalogue and includes additional assumptions; source publication does not establish first-ever priority or certify this exact model remains unsolved. The survey explicitly asks how much debt markets can absorb. This probabilistic safety-capacity definition is an adapted operational target, not a published universal fiscal limit.","statusReviewedAt":"2026-10-08"}
% FIRST_POSTED: 2026-10-08
% ENTRY_KIND: statement_only
% !TeX program = lualatex
\documentclass[11pt]{article}
\usepackage{fontspec}
\usepackage[a4paper,margin=25mm]{geometry}
\usepackage{amsmath,amssymb,mathtools}
\usepackage{xurl}
\usepackage[hidelinks]{hyperref}
\newcommand{\sref}[2]{\hyperlink{source:#1:#2}{[S#2]}}
\setlength{\emergencystretch}{3em}
\begin{document}
% problemId:30007535
\section*{Safe debt and fiscal limits}\label{30007535}
\subsection{Definitions and mathematical statement}
Fix a sovereign issuing debt in its own currency, a fiscal reaction rule, monetary institutions, and a stochastic process for tax capacity and expenditure commitments. Let \(b\) be the initial debt-to-output ratio. Debt delivers an explicitly modeled convenience service to investors, but government bonds may also face rollover failure, restructuring, or inflation-related real losses. For horizon \(H\), let \(L_H(b)\) be the maximum cumulative real loss to a holder under the induced equilibrium and define a safety event \(L_H(b)\le\ell\), with declared tolerance \(\ell\ge0\).
For risk tolerance \(\alpha\in(0,1)\), define a conditional safe-debt capacity
\[
b^{\mathrm{safe}}=
\sup\{b\ge0:\Pr(L_H(b)>\ell)\le\alpha
\text{ and the fiscal rule is feasible over }H\}.
\]
Take the supremum in the extended nonnegative reals, with value zero when the feasible set is empty. Identify or bound this capacity and the marginal convenience benefit of additional debt using explicitly justified restrictions on investor demand, fiscal adjustment, and equilibrium selection. Distinguish interest–growth arithmetic from feasible contingent primary surpluses and from self-fulfilling rollover equilibria. Report sensitivity to tail events and institutional credibility. The definition is intentionally horizon-, loss-, and institution-specific; there need not be a universal debt threshold. Determine whether proposed additional issuance raises social welfare before the safety constraint binds, accounting for liquidity benefits and fiscal or inflation risks.

\par\textbf{Formulation and scope.} The survey explicitly asks how much debt markets can absorb. This probabilistic safety-capacity definition is an adapted operational target, not a published universal fiscal limit.

\par\textbf{Open-question provenance.} An explicit question or research agenda was verified at the source locator. This does not by itself certify that every later version remains unresolved \sref{30007535}{1}.

\par\textbf{Historical priority.} This is the earliest explicit relevant statement verified in this audit, not a claim of original priority; older literature and earlier versions may contain antecedents.

\subsection{Short English statement}
Identify a horizon- and institution-specific safe-debt capacity satisfying an explicit bond-loss probability constraint and fiscal feasibility.

\subsection{Sources}
\begin{itemize}\raggedright
\item[\textnormal{[S1]}]\hypertarget{source:30007535:1}{} Kris James Mitchener, Christoph Trebesch. 2021. Sovereign Debt in the 21st Century. Section 10, “Concluding Remarks,” pp. 51–52.
\url{https://www.ifo.de/DocDL/cesifo1_wp8959.pdf}.
DOI: 10.3386/w28598.
{\small\textit{Earliest verified question/agenda reference; historical priority qualified below. The survey says many puzzles remain, including how much more sovereign debt markets can absorb and how official holdings, bailouts, panics, and sovereign–bank feedback affect future debt crises.}}
\end{itemize}
\end{document}
